\newif\ifglobal

%%%%%%%%%%%%%%%%%%%%%%%
%%% Compile to view %%%
%%%%%%%%%%%%%%%%%%%%%%%

%%%%%%%%%%%%%%%%%%%%%%%%%%%%%%%%%%%%%%%%%%%%%%%%%%%
%%% Readme:                                     %%%
%%% https://www.overleaf.com/read/bwdjvfjksvjy/ %%%
%%%%%%%%%%%%%%%%%%%%%%%%%%%%%%%%%%%%%%%%%%%%%%%%%%%

% >>> M?: marks a module setting number or important notice

% >>> M4: Set {./relative-path} to external files for this module <<<
\def \modulefiles {./25-RNG}
% To add an external file, use:
% (Replace '---' with actual filename)
% '\input{\modulefiles/tables/---.tex}' for tables
% '\includegraphics{\modulefiles/figures/---}' for figures
% '\subfile{\modulefiles/---}' for register files


% >>> M1: This module is in global project? <<<
% 'YES' - uncomment; 'NO' - comment out
\globaltrue

\ifglobal
    \documentclass[main\_\_EN.tex]{subfiles}

\else
    \documentclass[openany, 10pt]{book}

    % >>> M2: In ./00-shared/preamble-trm-module, also find
    %         settings for visibility of todo notes, watermark <<<
    \usepackage{preamble-trm-module}

    % >>> M3: Temporary labels in English,
    %         you can add your own labels there <<<
    \myexternaldocument{./00-shared/temp-labels-en}

\fi %\ifglobal


\setcounter{secnumdepth}{4}


% >>> M5: If this module needs any updates in
% './00-shared/preamble-trm-module.sty'
% (include additional package, etc.), contact your module PM
% or create the file './00-shared/changelog.tex' make the
% necessary updates yourself and document them in this file <<<


\begin{document}


% Sets language into which pre-defined bits of text are translated
\selectlanguage{english}

% Do not delete this block
\ifglobal
% Add the following front matter if
% this module is outside of global project
\else
    \InputIfFileExists{readme}{}{}
    {\let\clearpage\relax \listoftodos}
    {\let\clearpage\relax \listofdonetodos}
    \tableofcontents
    %\listoftables
    %\listoffigures
\fi


%%%%%%%%%%%%%%%%%%%%%%%%%
%%% TRM Module Begins %%%
%%%%%%%%%%%%%%%%%%%%%%%%%

\hypertarget{rng}{}
\chapter{Random Number Generator (RNG)}\label{mod:rng}

\section{Introduction}

The \chipname{} contains a true random number generator, which generates 32-bit random numbers that can be used for cryptographical operations, among other things.

\section{Features}

The random number generator in \chipname{} generates true random numbers, meaning they are derived from a physical process rather than an algorithm. Each number within the specified range has an equal probability of occurrence.

\section{Functional Description}\label{sec:rng-func-descr}

Every 32-bit value that the system reads from the \hyperref[regdesc:LPPERIRNGSYNCDATAREG]{LPPERI\_RNG\_DATA\_SYNC\_REG} register of the random number generator is a true random number. These true random numbers are generated based on the \textbf{thermal noise} in the system, the \textbf{asynchronous clock mismatch}, and the \textbf{asynchronous counter}.

\begin{itemize}
    \item \textbf{Thermal noise} comes from the high-speed ADC, or SAR ADC, or both. Whenever the high-speed ADC or SAR ADC is enabled, bit streams will be generated and fed into the random number generator through an XOR logic gate as random seeds.
    \item RC\_FAST\_CLK is an \textbf{asynchronous clock source} source that can generate metastability in the circuit. This metastability can also be used as a random number seed \hyperref[others:rng-note-clock-seed]{$^2$} to feed into the random number generator.
    \item \textbf{Asynchronous counters} count using an asynchronous clock source. The asynchronous clock source generates circuit metastability, which can be used as a random number seed to feed into the random number generator. There are two types of asynchronous counters:
    \begin{itemize}
        \item On-chip \textbf{\nameref{mod:rtctimer} (RTC\_TIMER)}: The counter directly counts using the on-chip RTC timer. The count value is fed into the random number generator as a random number seed after applying XOR logic.
        \item \textbf{Ring oscillator (BUF\_CHAIN)}: The counter counts based on the ring oscillator, formed by connecting buffer cells end-to-end. The count value derived from this oscillator is fed into the random number generator as a random number seed after applying XOR logic.
    \end{itemize}
    These two types of asynchronous counters can work independently or in combination.
\end{itemize}

\begin{figure}[ht]
    \centering
    \includegraphics[width=0.8\textwidth]{\modulefiles/figures/25-RNG-noise-source}
    \caption{Noise Source}
    \label{fig:rng-noise-source}
\end{figure}

When there is noise coming from the SAR ADC, the random number generator is fed with a 2-bit entropy in one clock cycle (configurable, normally 66 kHz). Thus, it is advisable to read the \hyperref[regdesc:LPPERIRNGSYNCDATAREG]{LPPERI\_RNG\_DATA\_SYNC\_REG} register at a maximum rate of 66 kHz to obtain the maximum entropy.

When there is noise coming from the high-speed ADC, the random number generator is fed with a 2-bit entropy in one APB clock cycle, which is normally 80 MHz. Thus, it is advisable to read the \hyperref[regdesc:LPPERIRNGSYNCDATAREG]{LPPERI\_RNG\_DATA\_SYNC\_REG}
register at a maximum rate of 5 MHz to obtain the maximum entropy.

When BUF\_CHAIN is enabled, 1 bit of entropy can be obtained in each RC\_FAST\_CLK cycle.

A data sample of 2 GB, which is read from the random number generator at a rate of 5 MHz with only the high-speed ADC being enabled, has been tested using the Dieharder Random Number Testsuite (version 3.31.1). The sample passed all tests.


\section{Programming Procedure}

To use the random number generator, set \hyperref[fielddesc:LPPERIRNGCKEN]{LPPERI\_RNG\_CK\_EN} to enable its clock. The clock is automatically enabled when \nameref{mod:ecdsa} is used, regardless of
\hyperref[fielddesc:LPPERIRNGCKEN]{LPPERI\_RNG\_CK\_EN}.

When using the random number generator, make sure the entropy source is enabled. Otherwise, pseudo-random numbers will be returned.
\begin{itemize}
    \item SAR ADC can be enabled by using the DIG ADC controller. For details, please refer to Chapter \ref{mod:adc} \textit{\nameref{mod:adc}}.
    \item High-speed ADC is enabled automatically when the Wi-Fi or Bluetooth module is enabled.
    \item RC\_FAST\_CLK is always enabled
    \item Asynchronous counters
        \begin{itemize}
            \item RTC timer (RTC\_TIMER) is enabled by setting the \hyperref[fielddesc:LPPERIRTCTIMEREN]{LPPERI\_RTC\_TIMER\_EN} field.
            \item Ring oscillator (BUF\_CHAIN) is enabled by setting the \hyperref[fielddesc:LPPERIRNGSAMPLEENABLE]{LPPERI\_RNG\_SAMPLE\_ENABLE} bit.
        \end{itemize}
\end{itemize}
\newpage
\begin{tiplisting}\label{others:rng-note-clock-seed}
\vspace{-2.5em}
\begin{enumerate}
    \item When the Wi-Fi module is enabled, the high-speed ADC may produce saturated readings in extreme cases, which reduces entropy. To mitigate this, enable the SAR ADC as an additional noise source for the random number generator.
    \item Enabling RC\_FAST\_CLK and asynchronous counters increases the RNG entropy. However, to ensure maximum entropy, always enable the SAR ADC or high-speed ADC as well.
\end{enumerate}
\end{tiplisting}

When multiple random numbers are required, the value of the \hyperref[regdesc:LPPERIRNGSYNCDATAREG]{LPPERI\_RNG\_DATA\_SYNC\_REG} register can be read continuously. However, when reading the register, ensure that the rate does not exceed the entropy generation rate described in Section \ref{sec:rng-func-descr}. If multiple entropy sources are configured, it is recommended to not exceed the entropy generation rate of the lowest noise source.


%%%%%%%%%%%%%%%%%%%%%%%%
%%% Register Summary %%%
%%%%%%%%%%%%%%%%%%%%%%%%

\newpage
\hypertarget{rng-reg-summ}{}
\section{Register Summary}
\label{sec:rng-reg-summ}

The addresses in this section are relative to Random Number Generator base address provided in Table \ref{tab:sysmem-base-address} in Chapter \ref{mod:sysmem} \textit{\nameref{mod:sysmem}}.

The abbreviations given in Column \textbf{Access} are explained in Section \hyperref[glossary-access-types]{\textit{Access Types for Registers}}.

\begin{longtable}{ | p{6cm} | p{7.2cm} | p{1.3cm} | p{1cm} | }
\hline\rowcolor{lightgray}
\textbf{Name} & \textbf{Description} &  \textbf{Address} & \textbf{Access} \\ \hline
\endhead
\hyperref[regdesc:LPPERIRNGCFGREG]{LPPERI\_RNG\_CFG\_REG} & RNG configuration register & 0x0024 & varies \\ \hline
\hyperref[regdesc:LPPERIRNGSYNCDATAREG]{LPPERI\_RNG\_DATA\_SYNC\_REG} & RNG result synchronization register & 0x0028 & RO \\ \hline
\end{longtable}

%%%%%%%%%%%%%%%%%
%%% Registers %%%
%%%%%%%%%%%%%%%%%

\newpage
\section{Registers}

The addresses in this section are relative to Random Number Generator base address provided in Table \ref{tab:sysmem-base-address} in Chapter \ref{mod:sysmem} \textit{\nameref{mod:sysmem}}.

For how to program reserved fields, please refer to Section \hyperref[programming-reserved-field]{\textit{Programming Reserved Register Field}}.

\begin{register}{H}{LPPERI\_RNG\_CFG\_REG}{0x0024}\label{regdesc:LPPERIRNGCFGREG}
\regfield{(reserved)}{20}{12}{00000000000000000000}%
\regfield{LPPERI\_RTC\_TIMER\_EN}{2}{10}{11}%
\regfield{(reserved)}{9}{1}{000000000}%
\regfield{LPPERI\_RNG\_SAMPLE\_ENABLE}{1}{0}{{0}}%
\reglabel{Reset}\regnewline%
\begin{regdesc}\begin{reglist}
\label{fielddesc:LPPERIRNGSAMPLEENABLE}\item [LPPERI\_RNG\_SAMPLE\_ENABLE] Configures whether to enable BUF\_CHAIN.\\0: Disable\\ 1: Enable\\  (R/W)
\label{fielddesc:LPPERIRTCTIMEREN}\item [LPPERI\_RTC\_TIMER\_EN] Configures whether to enable the RTC timer before CRC.\\
0: Disable\\
1: Enable\\ (R/W)
\label{fielddesc:LPPERIRNGSAMPLECNT}\item [LPPERI\_RNG\_SAMPLE\_CNT] Represents the count value of BUF\_CHAIN. (RO)
\end{reglist}\end{regdesc}
\end{register}


\begin{register}{H}{LPPERI\_RNG\_DATA\_SYNC\_REG}{0x0028}\label{regdesc:LPPERIRNGSYNCDATAREG}
\regfield{LPPERI\_RND\_SYNC\_DATA}{32}{0}{00000000000000000000000000000000}%
\reglabel{Reset}\regnewline%
\begin{regdesc}\begin{reglist}
\label{fielddesc:LPPERIRNDSYNCDATA}\item [LPPERI\_RND\_SYNC\_DATA] Represents the RNG synchronization result. (RO)
\end{reglist}\end{regdesc}
\end{register}

\end{document}
